\documentclass[11pt]{article}

\usepackage[margin=1in]{geometry}
\usepackage[T1]{fontenc}
\usepackage{lmodern}
\usepackage{amsmath,amssymb,mathtools,bm}
\usepackage{booktabs}
\usepackage{enumitem}
\usepackage{xcolor}
\usepackage{graphicx}
\usepackage{microtype}
\usepackage[hidelinks]{hyperref}

\definecolor{navy}{RGB}{25,55,92}
\definecolor{muted}{RGB}{90,96,105}
\newcommand{\R}{\mathbb{R}}
\newcommand{\E}{\mathbb{E}}
\newcommand{\Prob}{\mathbb{P}}
\newcommand{\Ind}{\mathbf{1}}
\newcommand{\Dip}{\operatorname{Dip}}
\newcommand{\why}[1]{\par\noindent\textcolor{muted}{\small\itshape Why: #1}}

\hypersetup{
  pdftitle={Auditing Hidden Routing Regimes: Consolidated Method Note (v3)},
  pdfsubject={Consolidated algorithm note with controlled evaluation},
}

\title{\textcolor{navy}{Auditing Hidden Routing Regimes in Black-Box Models\\
\large Consolidated method note (v3): definition, inputs/outputs, assumptions,
algorithm,\\ and a controlled evaluation with known regimes}}
\author{Robert Romani}
\date{July 29, 2026}

\begin{document}
\maketitle

\begin{abstract}
This note consolidates the v2 algorithm note (\emph{routing\_audit\_v2}, July 24) and its
derivation guide into a single description of the \emph{current} method, incorporating the
experimental findings of July 28 that changed the algorithm substantially: modality (the dip
test), not dispersion, is the primary detection statistic; the hierarchical GP null generator
is no longer required; the probe family is decisive rather than a reporting stratification;
query allocation is a first-class design parameter; and local $K{=}2$ recovery of the routing
penalty is solved while general-$K$ global synchronization remains open. The note gives, in
order: the definition of a regime and the recoverable estimand
(\S\ref{sec:regime}); the method's inputs and outputs (\S\ref{sec:io}); the complete
assumption set, including two new preconditions surfaced by experiment
(\S\ref{sec:assumptions}); the complete algorithmic steps (\S\ref{sec:algorithm}); the
evidence that forced each change from v2 (\S\ref{sec:changes}); a new controlled simulation
with known true regimes evaluating recovery under a $40\times$ noise sweep and near regime
boundaries (\S\ref{sec:sim}); and a like-for-like comparison against the previous graph-based
(SBM) approach and three simple baselines (\S\ref{sec:baselines}). Status and open problems
close the note (\S\ref{sec:open}).
\end{abstract}

\section{Setting and the definition of a regime}
\label{sec:regime}

A deployed black-box system is modeled as
\[
  f(x)=f_{g(x)}(x)+\eta,
\]
where $g:\mathcal X\to[K]$ is an unobserved routing rule, $f_k$ is the response mechanism of
branch $k$, and $\eta$ is observation noise. \textbf{A regime is a pair}: a branch $f_k$
together with the region $g^{-1}(k)$ on which the router selects it. The auditor observes
only queries and responses through a response map $\psi$; in the present applications the
response is scalar ($V=1$), which collapses the vector machinery of v2 (whitening is division
by a scalar; the dip is applied to the responses directly, with no projection choice).

\paragraph{Identifiability ceiling.} Regimes are recoverable only up to behavioural
equivalence: two branches that answer identically under $\psi$ at the probed inputs are
indistinguishable in principle, and a sharp-but-continuous transition narrower than the
probed scales is observationally equivalent to a hard gate. All claims are therefore relative
to $\psi$, the probe design, and the declared smoothness class. (The controlled evaluation in
\S\ref{sec:sim} sharpens this: the dip statistic responds to \emph{jumps} and not to
continuous kinks, so the observational-equivalence class is narrower than v2 assumed.)

\paragraph{The estimand is the boundary, not the partition.} Detection has a radius: a gate
is visible only from anchors whose probe neighbourhood actually straddles it. The natural
estimand is therefore the boundary --- the level set where $g$ flips --- together with
boundary-adjacent territory; per-anchor regime classification is a derived quantity, and
``unflagged'' is never read as ``interior.'' The corrected ground-truth criterion used
throughout: an anchor is \emph{detectable} at scale $\sigma$ when the minority probe mass
across the gate, $\pi_{\mathrm{true}}(\sigma)$, is at least $0.05$; truth is mixing fraction,
not distance.

\section{Inputs and outputs}
\label{sec:io}

\paragraph{Inputs.}
(i) query access to $f$ (no labels, no internals, no gradients);
(ii) a set of anchors $x_1,\dots,x_n$ (real instances);
(iii) a probe family and a scale ladder $\sigma_1<\dots<\sigma_T$ --- \emph{on-manifold}
(kernel-weighted resampling of real inputs) as primary where the covariate distribution is
locally smooth (assumption A12), ambient Gaussian otherwise or as robustness check;
(iv) a response map $\psi$;
(v) a per-anchor query budget $m$ per scale, with the allocation rule of
\S\ref{sec:algorithm}.
No protected attribute enters the audit; it is used only in Stage~C scoring.

\paragraph{Outputs.}
Per anchor: a flag with a $p$-value (dip, Bonferroni over the scales actually dipped); the
recovered penalty $\widehat\Delta$ (component-mean gap of a $K{=}2$ mixture at the most
significant scale) with a consistency measure across scales; the soft local composition
$\widehat\pi$; the anchor orientation $\widehat z$ (which branch owns the unperturbed
response); and an \emph{abstention} when support or signal is insufficient --- a first-class
output, distinct from ``null-consistent.''
Audit level: an alignment claim --- do flags concentrate near candidate gate boundaries? ---
tested by spatial (or graph) randomization with radius-matched statistics; abstention and
negative-control rates; and, when per-anchor confirmatory claims are needed,
Benjamini--Yekutieli across anchors as a dependence-robust sensitivity.

\section{Assumptions}
\label{sec:assumptions}

B1--B10 carry over from v2 \S8 unchanged in content; stated compactly with diagnostics:

\begin{itemize}[leftmargin=1.8em,itemsep=1pt]
  \item \textbf{B1 stationarity.} One routing law throughout the audit (no stateful
  adversary). \emph{Diag:} interleaved disguised queries; batch drift tests.
  \item \textbf{B2 probe accessibility.} Recovered regimes are those some probe reaches;
  $\widehat K$ counts observable regimes. \emph{Diag:} cross-support comparison.
  \item \textbf{B3 local distinguishability.} Locally present regimes differ in response law
  with bounded local classification error. \emph{Diag:} held-out likelihood, stability.
  \item \textbf{B4 prevalence.} Minor components need soft mass $m\widehat\pi\ge n_{\min}$
  ($n_{\min}=20$ here). \emph{Diag:} report soft masses; merge or omit failures.
  \item \textbf{B5 within-regime continuity.} Branch response laws drift Lipschitz-slowly
  between neighbouring anchors. \emph{Diag:} centroid smoothness along the anchor graph.
  \item \textbf{B6 matching margin} and \textbf{B7 overlap connectivity} --- required only
  for general-$K$ synchronization (open; \S\ref{sec:open}); trivial at $K{=}2$, where
  components are ordered by mean.
  \item \textbf{B8 anchor fidelity.} The unperturbed query is routed by the same mechanism
  (used only for orientation). \emph{Diag:} repeated unperturbed queries.
  \item \textbf{B9 gate regularity} and \textbf{B10 transversality.} Locally linear boundary
  at probe scale; probes must have variance along the boundary normal. \emph{Diag:} held-out
  label prediction; support comparison.
\end{itemize}

Two preconditions the July experiments added, both \emph{checkable}:

\begin{itemize}[leftmargin=1.8em,itemsep=2pt]
  \item \textbf{A11 (honest-model smoothness at every probed scale).} The single-branch
  surface must be effectively straight at all ladder scales: lengthscale $\ell>\sigma_T$ and
  smooth amplitude far below the gap of interest. Verified for Seattle two independent ways:
  $\ell\approx 2367$\,m against a ladder top of $1200$\,m, and local curvature amplitude
  $s\approx 0.004$ log-EUI against a penalty of $0.30$ ($75\times$). The simulation
  (\S\ref{sec:sim}) shows the condition is \emph{sharp}: GP nulls with $\ell$ above the
  ladder top flag at $0.00$--$0.04$; with $\ell$ below it, at $0.33$--$0.81$. Operationally:
  \emph{cap the ladder below the honest lengthscale}, or declare the smoothness class ---
  this is the surviving content of v2's $\ell_{\min}$ registration, now attached to the dip
  rather than to a null generator.
  \item \textbf{A12 (covariate local smoothness, for on-manifold probes).} The probe
  distribution itself must be locally smooth in the probe metric; otherwise the data's own
  clustering is read as model multimodality. This is what defeated the attribute-graph
  experiment (construction-boom years, standard floor counts: dip $p\approx 0$ on the raw
  covariates). Checkable by dipping the covariates; ambient probes satisfy it by
  construction, which partially offsets their off-manifold disadvantage. Geography satisfied
  it by luck, not principle.
\end{itemize}

\section{The algorithm}
\label{sec:algorithm}

The simplified pipeline. Stage A$'$ makes the existence and localization claims; Stage B$'$
recovers local structure at $K{=}2$; Stage C$'$ makes the audit-level claims. The dispersion
apparatus of v2 \S3.1--3.4 and the hierarchical GP null generator of v2 \S3.3 are \emph{not
part of the pipeline} (see \S\ref{sec:changes} for the evidence).

\paragraph{Stage A$'$: detection by naive multi-scale dip scan.}
\begin{enumerate}[label=\textbf{A\arabic*.},leftmargin=2.7em,itemsep=2pt]
  \item \textbf{Register} the design: $\psi$; probe families; a ladder of $3$--$6$
  log-spaced scales capped by A11; per-scale budget $m$; $\alpha$; the effect-size
  threshold; negative controls; support diagnostics.
  \item \textbf{Probe} each anchor at each scale: on-manifold primary (kernel-weighted
  resample of real inputs, weights $\propto e^{-d^2/2\sigma^2}$), abstaining at scales with
  fewer than $12$ distinct neighbours; ambient as robustness family.
  \item \textbf{Minimum-signal rule, per scale:} dip a scale only if the response variance
  clears the observation-noise floor ($\chi^2$ test at $5\%$). Scales failing it are
  abstained from, not counted as evidence of absence. \why{applied per scale, not per
  anchor; the first Seattle run produced $R^{\log}\approx 18$ purely from an epsilon clamp
  when this was done per anchor.}
  \item \textbf{Dip} the (whitened) responses at each surviving scale; \textbf{flag} the
  anchor if $\min_t p_t < \alpha / \#\{\text{scales dipped}\}$ (Bonferroni across scales).
  No dispersion curve, no $\sigma^\star$.
  \item \textbf{Query allocation.} Power is limited by the minor-mode mass $m\,\pi$, not by
  the test. Prefer fewer scales with larger $m$: on the weak band
  $\pi_{\mathrm{true}}\in[0.05,0.10]$ at equal budget, $3\times m{=}1000$ gives power
  $0.875$ versus $0.472$ for $6\times m{=}500$. Design rule: keep $m\pi$ above the dip's
  power threshold in the weakest band that matters.
\end{enumerate}

\paragraph{Stage B$'$: local $K{=}2$ recovery.}
\begin{enumerate}[label=\textbf{B\arabic*.},leftmargin=2.7em,itemsep=2pt]
  \item \textbf{Fresh draw} per scale (budget split respected: the recovery sample is never
  the detection sample).
  \item \textbf{Mixture.} Fit a $2$-component Gaussian mixture; order components by mean;
  retain soft responsibilities. Report $\widehat\pi$ (mean responsibility), and
  $\widehat\Delta$ = component-mean gap \emph{at the scale of minimum dip $p$-value} (the
  across-scale median dilutes $\widehat\Delta$ with non-crossing scales), plus its
  cross-scale consistency (CV). Apply the minimum-mass rule $m\widehat\pi\ge 20$.
  \item \textbf{Effect-size filter.} Keep flags with $\widehat\Delta$ above a registered
  threshold (Seattle: $\widehat\Delta>0.15$ lifted precision $0.968\to 1.000$ at $652$
  flags). Report as a precision/recall curve, not a single point: $\widehat\Delta$ is biased
  downward at small $\pi$, so a fixed threshold trades recall in the weak band.
  \why{$\widehat\Delta$ and the dip read the same bimodality --- this is one test plus a
  practical-significance filter, not independent corroboration.}
  \item \textbf{Orientation.} Posterior-assign the unperturbed response to a component;
  at $K{=}2$ synchronization is ordering by mean. General $K$: open
  (v2 SYNC1--7 remains the committed procedure; see \S\ref{sec:open}).
\end{enumerate}

\paragraph{Stage C$'$: audit-level claims.}
\begin{enumerate}[label=\textbf{C\arabic*.},leftmargin=2.7em,itemsep=2pt]
  \item \textbf{Alignment by randomization.} Rather than $3{,}581$ dependent per-anchor
  claims, test what the audit cares about: do flags concentrate near candidate boundaries?
  Permute the gating covariate across units (tracts), holding geometry and anchors fixed
  --- valid under arbitrary spatial dependence, no null generator. \textbf{Radius-match} the
  test statistic to the method's measured reach and report the full radius sweep
  (Seattle, on-manifold, $\le 800$\,m: $p=0.0025$; a mismatched $\le 100$\,m statistic
  penalizes the method's own detection radius).
  \item \textbf{Per-anchor claims} (if needed): BY at $q=0.10$ across anchors as the
  dependence-robust correction (Seattle flags survive the $c(n)=8.68$ penalty).
  \item \textbf{Harm read-out} (optional): group-conditional conformal coverage. Seattle:
  marginal coverage $0.898$ at nominal $90\%$, but $0.840$ on the penalized branch versus
  $0.903$ off it --- invisible marginally, and the recovered partition is exactly the
  grouping a Mondrian conformal predictor needs.
  \item \textbf{Report} existence, structure, interpretation separately, with abstentions,
  controls, and certified-versus-extrapolated territory as in v2.
\end{enumerate}

\section{What changed from v2, and the evidence that forced it}
\label{sec:changes}

\begin{itemize}[leftmargin=1.8em,itemsep=2pt]
  \item \textbf{Dispersion demoted from primary statistic to (at most) optional
  corroborator.} It fails all three roles tested: not a detector (uncalibrated $R^{\log}$
  flags smooth GP nulls at $0.65$--$1.00$; calibrated $T^{\mathrm{scale}}$ reaches power
  $0.025$ or $0.000$ depending on the null's construction); not a scale selector
  ($\sigma^\star$ is \emph{worse than a random scale}: power $0.715$ vs $0.798$, because
  $r(\sigma)$ peaks small while the dip wants balanced mixing); not an efficient screen (AUC
  $0.750$; retaining $90\%$ of flags still requires dipping $70\%$ of anchors, at
  $\sim 4\times$ the queries it saves). Its surviving corroboration value (precision
  $0.52\to 0.81$, $\kappa=0.32$) is dominated by the $\widehat\Delta$ filter.
  \item \textbf{Modality promoted to primary.} A hard gate seen through a probe of scale
  $\sigma$ is observationally a transition of width $\approx\sigma$, and a GP with
  $\ell\approx\sigma$ reproduces the same dispersion curve --- \emph{no statistic computed
  from $r(\sigma)$ alone can separate them}. What differs is the response
  \emph{distribution}: a gate makes two clumps separated by the penalty. Seattle: dip flag
  rate $0.213$ gated vs $0.000$ honest across $3{,}420$ buildings.
  \item \textbf{Null generator dropped.} It existed to calibrate dispersion; with dispersion
  gone, the $\ell_{\min}$ registration problem, the contamination/cross-fitting machinery,
  and the misspecification exposure go with it. The surviving smoothness declaration is A11,
  attached directly to the ladder.
  \item \textbf{Probe family is decisive, not a caveat.} The spatial alignment claim is
  significant on-manifold ($p=0.013$--$0.033$ across all four statistics) and never ambient
  ($p=0.083$--$0.208$): long-range ambient flags dilute the alignment. On-manifold is
  primary \emph{where A12 holds} --- the attribute-graph negative result is precisely an A12
  violation, and ambient probes are the fallback there.
  \item \textbf{Estimand: boundary-adjacent territory} (detection radius $\approx$ probe
  reach; \S\ref{sec:sim} table by distance), with per-anchor labels derived.
  \item \textbf{Query allocation first-class} (A5 above): reallocation lifts projected
  Seattle power on detectable anchors from $0.737$ to ${\sim}0.94$.
  \item \textbf{Stage B split.} Local $K{=}2$ recovery with a pure anchor is \emph{solved}
  (Seattle: $\widehat\Delta=0.291$ against a planted $0.300$; honest-model placebo $0.014$).
  General-$K$ synchronization, certification, and the recovery theorem remain open.
\end{itemize}

\section{Controlled evaluation with known regimes}
\label{sec:sim}

New in this consolidation (code and outputs in \texttt{sim/}; settings registered in file
headers; seed $20260729$). Two constructions, both with the audit blind to ground truth:
\textbf{1-D}: $f(x)=\beta x-\Delta\Ind[x\ge 0.5]$, $\beta=0.15$, $\Delta=0.30$, ladder
$\{0.02,0.063,0.2\}$, $m_{\mathrm{dip}}=m_{\mathrm{rec}}=1000$; matched no-gate control;
$240$ anchors. \textbf{2-D}: wavy boundary $x_2=0.5+0.08\sin 4\pi x_1$, honest model linear
plus a long-lengthscale GP term ($\ell=0.5>\sigma_T=0.2$; amplitude $15\times$ below
$\Delta$), $480$ anchors --- the Seattle geometry class. Both satisfy the measured Seattle
ordering (A11), and both report the corrected ground truth $\pi_{\mathrm{true}}$.

\paragraph{Noise sweep (Peng's first axis).} $\tau_{\mathrm{obs}}$ swept $0.005\to 0.2$
($\Delta/\tau$ from $60\times$ to $1.5\times$):

\begin{center}\small
\begin{tabular}{lcccccc}
\toprule
$\tau_{\mathrm{obs}}$ & 0.005 & 0.01 & 0.02 & 0.05 & 0.10 & 0.20\\
\midrule
power (detectable anchors)      & 0.976 & 0.964 & 0.892 & 0.567 & 0.012 & 0.000\\
FP, matched no-gate model       & 0.000 & 0.000 & 0.000 & 0.000 & 0.000 & 0.004\\
abstention, no-gate model       & 0.000 & 0.000 & 0.000 & 0.000 & 0.325 & 0.792\\
$\widehat\Delta$ at $p_{\min}$ scale (true $0.30$) & 0.248 & 0.248 & 0.248 & 0.252 & 0.288 & ---\\
\bottomrule
\end{tabular}
\end{center}

False positives stay at or below $0.004$ across the entire sweep; when noise drowns the
signal the method \emph{abstains} rather than flags (the per-scale minimum-signal rule); and
wherever anything is flagged, the recovered penalty is within $0.05$ of truth with a mild,
explainable downward bias (the mixture absorbs trend width into the components).

\paragraph{Near regime boundaries (Peng's second axis).} Power against the true mixing
fraction, the controlled analogue of the Seattle power-vs-$\pi$ curve:

\begin{center}\small
\begin{tabular}{lcccccc}
\toprule
$\pi_{\mathrm{true}}$ & $\le .01$ & $.01\!-\!.05$ & $.05\!-\!.10$ & $.10\!-\!.20$ & $.20\!-\!.35$ & $.35\!-\!.50$\\
\midrule
flag rate, 1-D & 0.00 & 0.00 & 0.55 & 1.00 & 1.00 & 1.00\\
flag rate, 2-D & 0.00 & 0.00 & 0.71 & 1.00 & 1.00 & 1.00\\
\bottomrule
\end{tabular}
\end{center}

Detection is binary in $\pi$: perfect above $0.10$, impossible below $0.05$ ($m\pi<50$
minor-mode points), partial in between --- confirming that power is limited by minor-mode
mass, not by the test, and motivating the allocation rule. By distance (2-D): flag rate
$1.00$ out to the ladder top, $0.73$ at $1$--$1.5\times$, $0.015$ beyond --- the detection
radius \emph{is} probe reach. Orientation accuracy on flagged anchors: $0.956$; overall 2-D:
power $0.929$, FP $0.000$, $\widehat\Delta$ median $0.270$ ($0.297$ within $0.02$ of the
boundary).

\paragraph{Robustness, and a sharp new precondition (A11).} Smooth GP confounders and a
kink control at $\tau=0.02$:

\begin{center}\small
\begin{tabular}{lccc}
\toprule
null model & $\ell$ vs ladder top & FP rate & reading\\
\midrule
GP $A{=}.10,\ \ell{=}.30$ & above & 0.000 & precondition holds\\
GP $A{=}.20,\ \ell{=}.30$ & above & 0.042 & precondition holds\\
GP $A{=}.10,\ \ell{=}.10$ & straddles & 0.375 & marginal\\
GP $A{=}.10,\ \ell{=}.05$ & below & 0.325 & violated --- designed failure\\
GP $A{=}.20,\ \ell{=}.05$ & below & 0.808 & violated --- designed failure\\
kink (continuous, unrouted) & --- & 0.050 & \emph{not} flagged; dip targets jumps\\
\bottomrule
\end{tabular}
\end{center}

Dipping at scales beyond the honest lengthscale reads smooth wiggles as modes; capping the
ladder (checkable; Seattle had a factor of two to spare) restores nominal error. The kink
control refines v2's observational-equivalence discussion: the dip does \emph{not} flag
sharp-but-continuous behaviour at nominal rates --- modality is specific to jumps, a
specificity dispersion never had.

\begin{figure}[t]
\centering
\includegraphics[width=\textwidth]{fig_sim_known_regimes.png}
\caption{Known-regimes simulation. Left to right: noise sweep (power, FP, abstention);
power against true mixing fraction near the boundary (1-D and 2-D); penalty recovery
against noise with the honest-model placebo; the A11 precondition is sharp and checkable.}
\end{figure}

\begin{figure}[t]
\centering
\includegraphics[width=0.9\textwidth]{fig_sim2d_map.png}
\caption{2-D spatial gate: flags trace the wavy boundary (left); recovered
$\widehat\Delta$ per flagged anchor against a planted $0.30$ (right).}
\end{figure}

\section{Comparison with the graph-based approach and simple baselines}
\label{sec:baselines}

All methods run on the same 2-D gated/honest pair. The SBM row runs the April
behavioural-fingerprint thread's own code (\texttt{geospatial-xai-attacks/src/sbm}: kNN
cosine graph over per-anchor response fingerprints, collapsed-Gibbs Bernoulli SBM, $K{=}2$,
modal partition).

\begin{center}\small
\begin{tabular}{lccccc}
\toprule
method & existence FP & partition acc. & partition acc. & $\widehat\Delta$ err & queries\\
       & (honest)     & (all anchors)  & (detectable)   &                      & /anchor\\
\midrule
dip-scan (ours)        & 0.000 & \textbf{0.80} & \textbf{0.93} & \textbf{0.03} & 6001\\
global cluster + BIC   & 0.00  & 0.54 & 0.54 & 0.30 & 25\\
residual outlier       & 0.008 & 0.52 & 0.52 & --- & 25\\
SPA profile-mean       & \emph{no test} & 0.52 & 0.58 & 0.21 & 50\\
plain SBM, $K{=}2$ (their code) & \emph{no test} & 0.61 & --- & --- & 33\\
\bottomrule
\end{tabular}
\end{center}

Readings. \textbf{Global clustering} answers \emph{existence} in this easy setting (BIC
picks $K{=}2$ gated, $K{=}1$ honest) but its partition is chance-level: the detrending fit
absorbs most of the gate, and the residual split follows covariates rather than mechanism ---
the failure v2 \S5 predicted. \textbf{Residual outliers} see almost nothing (the smooth fit
absorbs the level shift). \textbf{SPA} has no calibrated existence claim and near-chance
partition; its vertex gap over-estimates from noise extremes ($\widehat\Delta$ error $0.21$)
--- averaging destroys the modality evidence, which was the original reason for abandoning
the route. \textbf{The SBM} aligns only weakly with the true regime (ARI $0.048$ gated,
$0.000$ honest) and returns a partition unconditionally: behavioural-fingerprint similarity
cancels the constant penalty (a level effect) except across boundaries, so the block
structure has little to latch onto --- intrinsic to similarity-graph methods, not a tuning
failure. Ours is the only method with a calibrated per-anchor existence claim, boundary
localization, and penalty recovery; its cost, $\sim 200\times$ the query budget of the cheap
baselines, is the price of those guarantees and is the axis the allocation rule optimizes.

\section{Status and open problems}
\label{sec:open}

\textbf{Demonstrated} (Seattle, $n=3{,}581$, planted gate $0.30$ on tract share, and the
controlled suite above): existence with $0.0075$ honest flag rate and power $0.737$
($\to\sim 0.94$ reallocated); precision $1.000$ at the registered operating point
($\widehat\Delta>0.15$, $652$ flags); recovery $\widehat\Delta=0.291$
(IQR $[0.152,0.299]$) vs placebo $0.014$; alignment $p=0.0025$ radius-matched;
conditional-coverage harm invisible to a marginal conformal check.

\textbf{Open, in rough order of value:}
\begin{enumerate}[leftmargin=2.2em,itemsep=1pt]
  \item \emph{General-$K$ synchronization and certification} (v2 SYNC1--7, CERT1--4) with
  the theorem-shaped target of v2 \S13: forest propagation correct with probability
  $\ge 1-\sum_e\alpha_e$ under B3--B7; unblocked but unstarted.
  \item \emph{Generalization beyond geospatial}: the anchor graph need not be geographic
  (v2 SYNC1), but A12 must hold in the probe metric. Routes: discrete graphs (feeder, rate
  class, tariff), proper residualization (the null generator's job resurfacing in a smaller
  role --- current partial fix reduces honest FP $0.27\to 0.10$, still $2\times$ nominal),
  or a smoother dataset.
  \item \emph{$\widehat\Delta$ at small $\pi$}: downward bias ($0.158$ on the Seattle weak
  band); report filter as precision/recall curve; pooling across the anchor graph is the
  likely fix and is the same machinery as (1).
  \item \emph{Adaptive adversaries} (response smoothing near probe scales) --- scoped in v2,
  untouched.
\end{enumerate}

\paragraph{Reproducibility.} The simulation suite is self-contained (numpy/pandas/
matplotlib; optional scipy/diptest/sklearn), seeded, and resumable; the Seattle pipeline's
reorganization into \texttt{geospatial-xai-attacks} with configs and a single
\texttt{run\_all} entry point accompanies this note.

\end{document}
